\chapter{Introduction}

This document contains my notes from Swansea University's graduate entry medicine course.

\section{Organisation}

The main body of the handbook is divided into chapters. Each chapter is
stored as a separate file inside the \texttt{chapters} folder.

For example:

\begin{itemize}
    \item \texttt{chapters/02-cardiovascular.tex}
    \item \texttt{chapters/03-respiratory.tex}
    \item \texttt{chapters/04-gastrointestinal.tex}
\end{itemize}

The chapters are imported in the required order from
\texttt{main.tex}. Keeping one topic in each file makes the project
easier to edit and reduces the chance of merge conflicts when several
people contribute.

\section{Suggested chapter structure}

A clinical topic might be organised under the following headings:

\begin{enumerate}
    \item definition and overview;
    \item relevant anatomy and physiology;
    \item aetiology and risk factors;
    \item pathophysiology;
    \item presentation;
    \item investigations;
    \item management;
    \item complications;
    \item prognosis and follow-up; and
    \item key revision points.
\end{enumerate}

Not every topic will require every heading.

\section{Formatting conventions}

Several reusable boxes are provided by \texttt{formatting.sty}.

\begin{keypoint}
Use this box for important principles or material that should be
remembered.
\end{keypoint}

\begin{clinicalnote}
Use this box for clinical context, practical interpretation, or links
between scientific knowledge and clinical practice.
\end{clinicalnote}

\begin{redflag}
Use this box for potentially serious presentations, urgent
considerations, or common safety errors.
\end{redflag}

\begin{examtip}
Use this box for assessment advice, useful contrasts, and easily
confused concepts.
\end{examtip}

\begin{definitionbox}[Example definition]
Use this box when introducing a particularly important term or
classification.
\end{definitionbox}

Important terms can also be written with
\verb|\term{important term}|, producing \term{important term}.
Generic drug names can be written with
\verb|\drug{drug name}|, producing \drug{drug name}.

\section{Sources and maintenance}

Where possible, notes should identify the source and date of clinical
recommendations. Guidance, diagnostic criteria, drug information, and
reference ranges can change over time.

Each chapter can end with a last-updated marker:

\begin{verbatim}
\lastupdated{2 September 2026}
\end{verbatim}

\begin{redflag}[Appropriate use]
These notes must not contain patient-identifiable information. They are
not a substitute for local protocols, current clinical guidelines,
formal prescribing resources, or advice from an appropriately
qualified professional.
\end{redflag}
